\begin{frame}{Evidence}

\begin{columns}[T,onlytextwidth]
\begin{column}{.48\textwidth}
\begin{block}{\faUserGraduate~Students \& graduates}
\begin{itemize}
    \item high engagement without attendance requirements
    \item transfer of reproducible workflows
    \item stronger documentation practices
\end{itemize}
\end{block}
\end{column}

\begin{column}{.48\textwidth}
\begin{block}{Academic and industry perspectives}
\begin{itemize}
    \item[\faBuildingColumns] more ambitious, student-driven projects
    \item[\faIndustry] improved coherence and credibility
\end{itemize}
\end{block}
\end{column}
\end{columns}

\end{frame}
\note{
There is no formal attendance requirement in this course.  
However, compared to other courses I teach with the same student group,  
attendance and participation are noticeably higher.  

In more conventional formats -- lectures, ``regular'' industry-facing final project -- engagement tends to be lower.  

Here, students are more actively involved  
throughout the process.  

We also observe transfer.  
Students adopt reproducible workflows and stronger documentation practices,  
and continue to use them beyond the course.  

After graduation, students often reflect on this format  
as one of the most impactful parts of their studies --  
not just in terms of content,  
but in how they learned to work.  

From the industry perspective,  
the outputs are more coherent and credible,  
and importantly, they can see how the results  
could be implemented in practice.  

So while this is not controlled evidence,  
the signal is consistent across students, graduates, and partners.
}